% Part: first-order-logic
% Chapter: proof-systems
% Section: introduction

\documentclass[../../../include/open-logic-section]{subfiles}

\begin{document}

\iftag{FOL}
      {\olfileid{fol}{prf}{int}}
      {\olfileid{pl}{prf}{int}}

\olsection{引言}

逻辑学通常兼具语义学与推导系统。
语义学涉及诸如真、可满足性、有效性和语义后承等概念。
推导系统的目的则是提供一种确立语义后承与有效性的纯句法方法。
之所以说它们是纯句法的，是因为在这样的系统中，一个推导就是一个有限的句法对象，通常是一个语句或公式的序列（或其他有限的排列方式）。
良好推导系统的性质是：任何给定的语句或公式序列（或排列），都可以被机械地验证为“正确”。

一阶逻辑最简单（也是历史上最早）的推导系统是\emph{公理化的}。
在这样的系统中，一个公式序列构成一个推导，当且仅当其中每个单独的公式要么属于一个固定的“公理”集，
要么通过有限多条“推理规则”从序列中位于它前面的公式推出——并且可以机械地验证一个公式是否为公理，
以及它是否通过某条推理规则从其他公式正确推出。
公理化推导系统易于描述——也易于在元理论层面处理——但其中的推导难以阅读和理解，也难以构造。

人们开发了其他推导系统，旨在便于构造推导，或使推导完成后更易理解。
例如自然演绎、真值树（也称为分析表证明）以及矢列演算。
有些推导系统特别为机械化而设计，例如归结法易于在软件中实现（但它的推导根本无法理解）。
这些其他推导系统大多将推导表示为公式的树，而非序列。
这使得便于看清推导的各个部分依赖于哪些其他部分。

因此，对于给定的逻辑（如一阶逻辑），不同的推导系统会对什么是语句为\emph{定理}、
以及语句从其他语句可推导意味着什么，给出不同的界定。
无论采用何种方式（公理推导、自然演绎、矢列推导、真值树、归结反驳），
我们都希望这些关系与语义上的有效性和语义后承概念相匹配。
我们用 $\Proves !A$ 表示“$!A$~是定理”，用 $\Gamma \Proves !A$ 表示“$!A$~可从~$\Gamma$ 推导”。
无论 $\Proves$~如何定义，我们都希望它与 $\Entails$ 匹配，即：

\begin{enumerate}
\item $\Proves !A$ 当且仅当 $\Entails !A$
\item $\Gamma \Proves !A$ 当且仅当 $\Gamma \Entails !A$
\end{enumerate}

上述条件的“仅当”方向称为\emph{可靠性}。
如果可推导性保证了语义后承（或有效性），则该推导系统是可靠的。
每个合格的推导系统都必须是可靠的；不可靠的推导系统完全无用。
毕竟，推导的全部目的就是为有效性或语义后承提供一种句法保证。我们将为所介绍的推导系统证明可靠性。

相反的“如果”方向也很重要：它称为\emph{完备性}。
一个完备的推导系统足够强大，能够证出只要 $!A$~有效，$!A$~就是定理；
并且只要 $\Gamma \Entails !A$，就有 $\Gamma \Proves !A$。
完备性更难确立，有些逻辑没有完备的推导系统。一阶逻辑则拥有完备的推导系统。
Kurt Gödel（哥德尔） 在其1929年的博士论文中首次证明了一阶逻辑推导系统的完备性。

另一个与推导系统相关的概念是\emph{一致性}。
如果从一个语句集合可以推导出任意语句，则称该集合为\emph{不一致的}；否则，称其为\emph{一致的}。
不一致性是不可满足性在句法上的对应物：就像不可满足的集合一样，
不一致的语句集合不能构成合格的理论，它们在根本层面上是有缺陷的。
虽然一致的语句集合未必为真或未必有用，但至少它们跨过了逻辑有用性的最低门槛。
对于不同的推导系统，语句集合一致性的具体定义可能有所不同；
但就像~$\Proves$ 一样，我们希望一致性能与其语义上的对应物——可满足性——相吻合。
我们希望始终满足：$\Gamma$ 是一致的，当且仅当它是可满足的。
这里，“仅当”方向相当于完备性（一致性保证可满足性），而“当”方向则相当于可靠性（可满足性保证一致性）。
事实上，对于经典一阶逻辑而言，可靠性与完备性的这两种表述是等价的。

\end{document}